\subsection{Preliminaries}

This section introduces existing models commonly used in similar studies.
They were selected to provide reference points for the case studies rather than to constitute a comprehensive model comparison.
Baseline models, such as linear interpolation and drift, provide simple references for interpreting more flexible and complex methods.
Statistical models, such as cubic spline interpolation and Theil–Sen regression, make assumptions about the relationships among variables and use data to estimate those relationships.
Machine learning models, such as ridge regression, are generally more flexible and can handle noisy data more effectively.
Together, these models allow the workflow to examine how conclusions vary across modeling assumptions.

\subsubsection{Baseline Models}

Baseline models are simple models used as references for evaluating the performance of more complex models.
In this study, we use linear interpolation and moving average interpolation as baseline models for time series imputation.
Linear interpolation estimates missing values by connecting the nearest known data points with a straight line.
Moving average interpolation estimates missing values by averaging known data points within a specified window around each missing value.

For time series forecasting tasks, we use naive persistence, drift, and Holt-style exponential smoothing~\cite{holt2004forecasting}.
Naive persistence assumes that the future value of a variable remains the same as its most recent observed value.
Drift extrapolates the trend observed in the historical data into the future.
Holt-style exponential smoothing captures both level and trend components in the data, allowing for more flexible forecasting.

\subsubsection{Statistical Models}

Statistical models use statistical methods to analyze data and make predictions.
In this study, we use cubic spline interpolation~\cite{de1978practical} for imputation tasks, which fits a piecewise cubic function to the data while ensuring smoothness at the known data points.
For forecasting tasks, we use local Theil--Sen regression~\cite{sen1968estimates}, a non-parametric method that estimates the slope of a trend by taking the median of all pairwise slopes.
We also include a non-seasonal Auto-ARIMA model~\cite{box2015time} as an additional statistical comparator.
Because the available time series are short, ARIMA results are interpreted only for the specific case studies analyzed here and are not taken as evidence of general performance.

\subsubsection{Machine Learning Models}

Machine learning models used in this study include polynomial ridge regression~\cite{hoerl1970ridge}, Gaussian process regression (GPR)~\cite{seeger2004gaussian}, and support vector regression (SVR)~\cite{smola2004tutorial}. We also use random forest~\cite{breiman2001random} for feature-importance analysis.

Polynomial ridge regression is a linear model that incorporates polynomial features and applies L2 regularization to prevent overfitting.

GPR is a Bayesian approach that models the underlying function as a Gaussian process. The trend is modeled as
\begin{equation}
	f(x) = \mathcal{GP}(0, k(x, x')),
\end{equation}
with an RBF-plus-noise kernel
\begin{equation}
	k(x, x') = \exp \left( -\frac{(x - x')^2}{2 \ell^2} \right) + \sigma^2 \delta_{x, x'},
\end{equation}
where $\ell$ controls temporal smoothness and $\sigma^2$ absorbs observation noise.

SVR is a regression method that finds a hyperplane in a high-dimensional space that best fits the data while allowing for some margin of error. The optimization problem for SVR can be formulated as
\begin{equation}
	\min_{w, b} \frac{1}{2} \| w \|^2 + C \sum_{i=1}^n \max(0, |y_i - (w^T \phi(x_i) + b)| - \epsilon),
\end{equation}
with the RBF kernel
\begin{equation}
	K(x_i, x_j) = \exp \left( -\gamma \| x_i - x_j \|^2 \right),
\end{equation}
where \(C\) controls the penalty for large residuals and \(\epsilon\) defines the width of the insensitive tube around the fitted function.

We also use isolation forest~\cite{liu2008isolation} as an anomaly detection method to identify outliers in the data.
It works by isolating observations that differ significantly from normal observations using random partitioning.
Because anomalies are few and distinct, they are isolated more quickly than normal data.
The contamination rate specifies the approximate fraction of observations that the method flags as potential anomalies.

Random forest~\cite{breiman2001random} is an ensemble learning method that trains many decision trees, each fitted to a bootstrap sample of the training data.
At each split, a random subset of features is used to reduce correlation among the trees.
The final prediction is obtained by aggregating the predictions of all trees.
In this study, random forest is used to analyze feature importance for a regression task, and for regression trees, node impurity is measured by the mean squared error of the target values within a node in \texttt{sklearn}.
At each split, the impurity decrease is calculated as
\begin{equation}
	\Delta I = I_P - \frac{n_L}{n}I_L - \frac{n_R}{n}I_R,
\end{equation}
where \(I_P\) is the impurity of the parent node, \(I_L\) and \(I_R\) are the impurities of the left and right child nodes, respectively, and \(n_L\) and \(n_R\) are the number of samples in the left and right child nodes, respectively, and \(n = n_L + n_R\).
The importance of a feature is then computed by summing the weighted impurity decreases over all splits where the feature is used, averaging this quantity across all trees in the forest, and normalizing the importances.

